\section{Caracterización de los datos y de la información requerida}

\subsection{Fuente y estructura}

Los datos provienen del archivo \emph{BDD\_Septiembre\_2026} de la ANT \cite{ref2}, formado por registros administrativos que levantan 16 entes de control en el lugar del siniestro. No son una muestra diseñada para estudiar la gravedad, y esa condición limita desde el inicio lo que puede concluirse. El archivo tiene 17.236 filas y 374 columnas, alrededor de 6,4 millones de celdas. De esas columnas, 301 describen a participantes ---siete campos para hasta 43 participantes por siniestro--- y las 73 restantes describen el evento y los vehículos; la columna \texttt{RANGO DE EDAD} aparece suelta dentro del bloque de participantes y no lleva sufijo. Como la mediana es de dos participantes por siniestro, las columnas de participantes están vacías en un 95,0\,\% y el 82,6\,\% de todas las celdas está vacío (Tabla~\ref{tab:a1}; Anexo~B).

El libro incluye además hojas auxiliares con los catálogos de causa probable, tipo de siniestro, tipo de vehículo y víctimas. Esas hojas funcionan como diccionario y fueron necesarias para interpretar la convención de fallecidos descrita en la sección~2.2; todas las cifras de esta sección se obtuvieron con el procedimiento del Anexo~B sobre la hoja \texttt{BDD\_2026}.

\subsection{Etiqueta, desbalance y variación en el tiempo}

Con la definición de la sección~2.2, 1.656 siniestros son graves (9,61\,\%). La mayoría registra un solo fallecido y solo 159 tienen dos o más (Tabla~\ref{tab:a1}). Esa clase minoritaria no está repartida de forma pareja en el tiempo: el número de graves por mes es casi constante, entre 172 y 196, pero el total de siniestros crece de 1.510 a 2.242. Como consecuencia, la proporción baja de 11,6\,\% en enero a 8,8\,\% en septiembre (Tabla~\ref{tab:a2}; Anexo~B).

Al mismo tiempo cambia la composición de los registros. Los partes de la Policía Nacional pasan de representar el 17,4\,\% del total en enero a un rango de 26,4\,\% a 28,9\,\% entre agosto y septiembre. La causa de ese cambio no está verificada y puede deberse a cobertura, a forma de registro o a un cambio real en la siniestralidad. Por eso se trata como un posible sesgo de registro o deriva temporal, no como una tendencia de seguridad vial. Por eso la evaluación respeta el orden temporal (sección~9).

\subsection{Calidad: valores faltantes y categorías desconocidas}

Las variables que describen el evento ---fecha, hora, provincia, cantón, zona, tipo de siniestro, causa y ente de control--- no tienen valores vacíos. Las limitaciones están en otras variables. La jerarquía y la ubicación de la vía valen \texttt{ND} en el 66,4\,\% de los registros; el 22,5\,\% no identifica el tipo del primer vehículo. Ambas aportan poca información y se tratan como categorías propias en lugar de imputarse.

Hay además defectos menores que exigen reglas explícitas (Tabla~\ref{tab:a3}; Anexo~B): diez registros escriben \texttt{colisión} en minúscula, cinco horas están almacenadas como texto, un mismo identificador aparece en dos siniestros de fechas distintas y catorce registros tienen valor~0 en la jerarquía de la vía. No hay filas duplicadas y las coordenadas no presentan vacíos ni valores fuera de un rango plausible, aunque no se usarán por privacidad. Ninguno de estos defectos cambia la etiqueta, pero todos deben quedar documentados para que otra persona obtenga las mismas cifras.

\subsection{Clasificación de las variables según disponibilidad y fuga}

La clasificación siguiente aplica las definiciones de las secciones~2.2 y~2.3 a las columnas del archivo.

\begin{table}[htbp]
\centering
\small
\caption{Variables según disponibilidad, fuga y privacidad}
\label{tab:variables}
\begin{tabularx}{\textwidth}{p{3.1cm}X p{4.6cm}}
\toprule
Grupo & Variables & Tratamiento \\
\midrule
Conjunto A & Fecha, hora y franja horaria, día, feriado, provincia, cantón, zona urbana o rural, zona de planificación, jerarquía de la vía, tipo de siniestro, tipo y servicio de los vehículos, número de vehículos, ente de control & Entrada. El ente de control puede reflejar cobertura (sección~6.6) y la jerarquía de la vía tiene 66,4\,\% de \texttt{ND} \\
Conjunto B & Código y descripción de la causa probable & Entrada solo en el escenario de contraste \\
Desenlace & \texttt{FALLECIDOS}, \texttt{LESIONADOS} & No son entrada: \texttt{FALLECIDOS} construye la etiqueta \\
Excluidas por fuga & \texttt{CONDICION\_1} a \texttt{CONDICION\_43} & Coinciden exactamente con el conteo de víctimas en los 17.236 registros (Tabla~\ref{tab:a4}); usarlas entregaría la respuesta \\
Excluidas por privacidad o ambigüedad & \texttt{EDAD}, \texttt{SEXO}, \texttt{RANGO DE EDAD}, \texttt{CASCO}, \texttt{CINTURON}, \texttt{TIPO\_ID}, \texttt{PARTICIPANTE}, coordenadas, dirección, nombre e identificador de la vía & No se usan, conforme a la restricción R-07 \\
\bottomrule
\end{tabularx}
\end{table}

La exclusión por fuga no es una suposición: el conteo de participantes en condición \texttt{FALLECIDO} coincide con \texttt{FALLECIDOS} en el 100\,\% de los registros, y lo mismo ocurre con \texttt{LESIONADO} y \texttt{LESIONADOS}. Esa evidencia sostiene la decisión D4.

\subsection{Variabilidad de la gravedad entre grupos}

La proporción de siniestros graves difiere lo suficiente entre grupos como para que la caracterización tenga sentido (Tabla~\ref{tab:a5}; Anexo~B). Según el tipo de siniestro va del 3,0\,\% en colisiones y caídas de pasajero al 45,8\,\% en arrollamientos. Es del 14,0\,\% en zona rural frente al 7,1\,\% en zona urbana, y del 13,4\,\% cuando interviene una motocicleta frente al 7,8\,\% cuando no. Entre entes de control, la Policía Nacional registra 16,3\,\% y la ATM de Guayaquil 6,7\,\%; entre causas probables con al menos 100 casos, el rango va del 1,9\,\% al 23,1\,\%. En la Tabla~\ref{tab:a5} se omiten los grupos con menos de 30 casos, de acuerdo con el criterio CA-03.

Estas diferencias son descriptivas y no implican causalidad. Varias de las variables que las producen están relacionadas entre sí ---el tipo de siniestro con el número de vehículos, por ejemplo---, y los grupos pequeños tienen proporciones inestables. Esa es la razón por la cual OE1 exige informar el número de casos y una medida de incertidumbre junto con cada proporción.

\subsection{Sesgos, dependencias y restricciones de uso}

El sesgo principal proviene del registro: el ente de control puede reflejar cobertura y no solo riesgo, de modo que las diferencias entre entes no deben leerse como diferencias entre territorios. Entre las variables hay dependencias que condicionan el análisis. La hora es cíclica, porque las 23:00 y las 00:00 son cercanas; los siniestros próximos en el tiempo y en el espacio comparten condiciones; y el número de vehículos se relaciona con el tipo de siniestro. Todo el análisis se apoya, además, en un solo periodo de nueve meses, sin comparación con otros años.

En cuanto al uso, los datos son públicos pero no se identificó una licencia explícita. Se usarán solo con fines académicos, citando a la ANT, y sin publicar registros individuales, conforme a la restricción R-07.

\subsection{Contraste con cifras oficiales}

Para comprobar que el archivo corresponde a la fuente oficial se contrastan sus totales. Para el primer trimestre de 2026 reúne 4.789 siniestros, 4.032 lesionados y 603 fallecidos, idénticos a los publicados por el INEC \cite{ref1}; para enero a mayo reúne 8.718 siniestros, igual al valor informado por la ANT para ese periodo \cite{ref7} (Tabla~\ref{tab:a4}; Anexo~B). La coincidencia es evidencia preliminar del criterio CA-01 y respalda que la etiqueta equivale a fallecidos en sitio.
